\chapter{Git basics}
\label{ch:git-basics}

\section{The idea}\label{the-idea}

The blocks labelled "What you should see" in this guide come from running these commands in a test folder. Your commit codes, dates and folder names will differ, \texttt{...} marks a shortened folder path, and very long output is trimmed.

Git is a program that records versions of a folder. Each saved version, called a \emph{commit}, stores what changed, who changed it, when and why. Nothing committed is lost, so you never need files such as \texttt{final\_v2\_really\_final.docx}.

Install it from \href{https://git-scm.com}{git-scm.com}\footnote{\texttt{https://git-scm.com}}, then tell it who you are, once:

\begin{lstlisting}[language=bash]
git config --global user.name "Your Name"
git config --global user.email "you@example.org"
\end{lstlisting}

\section{The everyday cycle}\label{the-everyday-cycle}

\begin{lstlisting}[language=bash]
git init -b main         # start tracking this folder, once, on a branch called main
git status               # what has changed?
git diff                 # show the changes, line by line
git add results.txt      # choose what goes in the next version
git commit -m "Correct the mean"   # save that version, with a reason
git log --oneline        # list the versions
\end{lstlisting}

\textbf{What you should see}

\begin{lstlisting}[language=text]
$ git init -b main
Initialized empty Git repository in .../project/.git/

$ git status
On branch main

No commits yet

Untracked files:
  (use "git add <file>..." to include in what will be committed)
	notes.txt
	results.txt

nothing added to commit but untracked files present (use "git add" to track)

$ git add results.txt
$ git commit -m "First results"
[main (root-commit) 4f2d1b0] First results
 1 file changed, 1 insertion(+)
 create mode 100644 results.txt
\end{lstlisting}

Now change the number in \texttt{results.txt} and look again. In \texttt{git\ status\ -\/-short}, \texttt{M} means modified and \texttt{??} means not tracked. In \texttt{git\ diff}, a line starting \texttt{-} was removed and a line starting \texttt{+} was added.

\textbf{What you should see}

\begin{lstlisting}[language=text]
$ git status --short
 M results.txt
?? notes.txt

$ git diff
diff --git a/results.txt b/results.txt
index 62b900e..de5c1f9 100644
--- a/results.txt
+++ b/results.txt
@@ -1 +1 @@
-The mean is 10.
+The mean is 12.

$ git add results.txt
$ git commit -m "Correct the mean"
[main 2f4cdbe] Correct the mean
 1 file changed, 1 insertion(+), 1 deletion(-)

$ git log --oneline
2f4cdbe Correct the mean
4f2d1b0 First results
\end{lstlisting}

Without \texttt{-b\ main}, this version of Git names the first branch \texttt{master} and prints a long hint about it. The rest of this guide uses \texttt{main}.

The \texttt{add} step exists so that one commit can hold one logical change, even if you changed other files too.

\section{Try it safely}\label{try-it-safely}

This script makes a new folder and goes through the whole cycle. Read it, then run it.

\textbf{practice.sh}

\lstinputlisting{../../../teaching/git-ci/examples/practice.sh}

\textbf{What you should see}

\begin{lstlisting}[language=text]
 results.txt | 2 +-
 1 file changed, 1 insertion(+), 1 deletion(-)
[main ab8a9f9] Revert "Add the median"
 Date: Sun Nov 1 16:26:34 2026 +0000
 1 file changed, 1 deletion(-)
ab8a9f9 Revert "Add the median"
bd9ec39 Add the median
6f32cdc Correct the mean
8c59d4e First results
The mean is 12.
\end{lstlisting}

The last line matters: the median has gone from the file, yet all four versions are still in the history.

\section{What to track}\label{what-to-track}

Track what you write by hand: text, code and small inputs. Don\textquotesingle t track what can be regenerated, large data, or anything secret. A file called \texttt{.gitignore} lists what to skip:

\begin{lstlisting}[language=text]
__pycache__/
*.log
build/
data/raw/
.env
\end{lstlisting}

With that file in place, and \texttt{run.log}, \texttt{.env} and \texttt{build/a.pdf} in the folder, Git no longer mentions them:

\textbf{What you should see}

\begin{lstlisting}[language=text]
$ git status --short
?? .gitignore
?? notes.txt
\end{lstlisting}

\section{Good commit messages}\label{good-commit-messages}

Write what changed and why, briefly, such as \texttt{Correct\ the\ living-wage\ year}. Commit often, in small steps, so that each can be understood or undone by itself.

